\subsection{Средние расстояния до фокуса}
\label{sec:mean-distance}
Что такое среднее расстояние планеты от Солнца? Вопрос не так прост, как кажется: ответ зависит от того, по какой переменной усреднять. Можно усреднять по времени, можно~--- по истинной аномалии, а можно~--- по эксцентрической. Среднее значение величины $f$ по аномалии $x$ за один оборот определим как
\begin{equation*}
    \langle f \rangle_x = \frac{1}{2\pi}\int\limits_0^{2\pi} f\,dx.
\end{equation*}

\paragraph{Связь дифференциалов аномалий} Продифференцируем уравнение Кеплера \eqref{eq:kepler-eq} и воспользуемся выражением \eqref{eq:kepler-eq-r-E} для расстояния до фокуса:
\begin{gather*}
    dM = (1 - e\cos E)\,dE = \frac{r}{a}\,dE,\\
    a\,dM = r\,dE.
\end{gather*}
Средняя аномалия по определению растёт равномерно, $M = 2\pi t / T$, поэтому $dt = T\,dM / (2\pi)$. По второму закону Кеплера радиус-вектор за время $dt$ заметает площадь $\frac{1}{2} r^2\,d\nu$, а секториальная скорость постоянна и равна $\pi a b / T$:
\begin{gather*}
    \frac{1}{2} r^2\,d\nu = \frac{\pi a b}{T}\,dt = \frac{\pi a b}{T} \cdot \frac{T\,dM}{2\pi} = \frac{ab}{2}\,dM,\\
    a\,dM = \frac{r^2}{b}\,d\nu = \frac{r^2}{a\sqrt{1 - e^2}}\,d\nu.
\end{gather*}
Соберём обе связи в одну цепочку:
\begin{equation}
    a\,dM = r\,dE = \frac{r^2}{a\sqrt{1 - e^2}}\,d\nu.
    \label{eq:anomaly-differentials}
\end{equation}
Из неё сразу следуют все три средних расстояния.

\paragraph{По истинной аномалии} Из правого равенства \eqref{eq:anomaly-differentials} $r\,d\nu = a\sqrt{1 - e^2}\,dE = b\,dE$, поэтому
\begin{equation*}
    \langle r \rangle_\nu
    = \frac{1}{2\pi}\int\limits_0^{2\pi} r\,d\nu
    = \frac{b}{2\pi}\int\limits_0^{2\pi} dE
    = b.
\end{equation*}

\paragraph{По эксцентрической аномалии} Из левого равенства $r\,dE = a\,dM$:
\begin{equation*}
    \langle r \rangle_E
    = \frac{1}{2\pi}\int\limits_0^{2\pi} r\,dE
    = \frac{a}{2\pi}\int\limits_0^{2\pi} dM
    = a.
\end{equation*}

\paragraph{По средней аномалии} Это среднее по времени. Из левого равенства $r\,dM = r^2\,dE / a = a (1 - e\cos E)^2\,dE$:
\begin{multline*}
    \langle r \rangle_M
    = \frac{1}{2\pi}\int\limits_0^{2\pi} r\,dM
    = \frac{a}{2\pi}\int\limits_0^{2\pi} \left(1 - 2e\cos E + e^2\cos^2 E\right) dE = \\
    = \frac{a}{2\pi}\left(2\pi - 0 + e^2 \pi\right)
    = a\left(1 + \frac{e^2}{2}\right).
\end{multline*}

Итак,
\begin{equation}
    \langle r \rangle_\nu = a\sqrt{1 - e^2},
    \qquad
    \langle r \rangle_E = a,
    \qquad
    \langle r \rangle_M = a\left(1 + \frac{e^2}{2}\right).
    \label{eq:mean-distances}
\end{equation}
\begin{wrapfigure}[15]{r}{0.5\tw}
    \centering
    \vspace{-0.5pc}
    \tikzsetnextfilename{mean-distance}
    \begin{tikzpicture}
        \footnotesize
        % Расстояние до фокуса как функция трёх аномалий, e = 0.7.
        % Все кривые параметризованы эксцентрической аномалией E
        % (в градусах): r = 1 - e cos E, M = E - e sin E,
        % tg(nu/2) = sqrt((1 + e) / (1 - e)) tg(E/2)
        \pgfmathsetmacro{\ecc}{0.7}
        \pgfmathsetmacro{\meanM}{1 + \ecc * \ecc / 2}
        \pgfmathsetmacro{\meanNu}{sqrt(1 - \ecc * \ecc)}
        \pgfmathsetmacro{\kNu}{sqrt((1 + \ecc) / (1 - \ecc))}
        \begin{axis}[
            width = 0.5\tw,
            height = 5.5cm,
            xmin = -180, xmax = 180,
            ymin = 0, ymax = 1.8,
            xtick = {-180, -90, 0, 90, 180},
            xticklabels = {$-\pi$, $-\pi/2$, $0$, $\pi/2$, $\pi$},
            xlabel = {аномалия},
            ylabel = {$r / a$},
            ytick = {0, 0.5, 1, 1.5},
            clip = false,
        ]
            \addplot [black, thick, domain = -180:180, samples = 91]
                ({x - \ecc * sin(x) * 180 / pi}, {1 - \ecc * cos(x)})
                node [pos = 0.78, anchor = south east, inner sep = 1pt] {$M$};
            \addplot [black, domain = -180:180, samples = 91]
                ({x}, {1 - \ecc * cos(x)})
                node [pos = 0.79, anchor = north west, inner sep = 1pt] {$E$};
            \addplot [black, thin, domain = -180:180, samples = 91]
                ({2 * atan(\kNu * tan(x / 2))}, {1 - \ecc * cos(x)})
                node [pos = 0.78, anchor = north west, inner sep = 1pt] {$\nu$};
            \addplot [thin, dashes] coordinates {(-180, \meanM) (180, \meanM)}
                node [pos = 1, anchor = west] {$\langle r \rangle_M$};
            \addplot [thin, dashes] coordinates {(-180, 1) (180, 1)}
                node [pos = 1, anchor = west] {$\langle r \rangle_E$};
            \addplot [thin, dashes] coordinates {(-180, \meanNu) (180, \meanNu)}
                node [pos = 1, anchor = west] {$\langle r \rangle_\nu$};
        \end{axis}
    \end{tikzpicture}
    \caption{Расстояние до фокуса как функция средней, эксцентрической и истинной аномалии, $e = 0.7$}
    \label{pic:mean-distance}
\end{wrapfigure}
Три средних упорядочены: $\langle r \rangle_\nu \leqslant \langle r \rangle_E \leqslant \langle r \rangle_M$, и совпадают они только для круговой орбиты, \lookPicRef{pic:mean-distance}. Для планет разница невелика~--- у Земли $e = 0.0167$, и все три средних отличаются от $a$ меньше чем на $0.02\%$. Для комет же разница огромна: у кометы Галлея $e = 0.967$ и $a = 17.8$~\au, откуда $\langle r \rangle_\nu = 4.5$~\au, а $\langle r \rangle_M = 26.1$~\au{} Комета почти всё время находится далеко за орбитой Сатурна, хотя около $80\%$ оборота по истинной аномалии проходит внутри орбиты Юпитера.

\paragraph{Среднее обратное расстояние} Тем же приёмом находится среднее по времени от $1/r$, которое определяет среднюю потенциальную энергию тела на орбите:
\begin{equation*}
    \left\langle \frac{1}{r} \right\rangle_M
    = \frac{1}{2\pi}\int\limits_0^{2\pi} \frac{dM}{r}
    = \frac{1}{2\pi}\int\limits_0^{2\pi} \frac{dE}{a}
    = \frac{1}{a}.
\end{equation*}
Средняя по времени потенциальная энергия равна $-GMm / a$, как у тела на круговой орбите радиуса $a$. Именно это соотношение используется при выводе полной энергии через теорему о вириале, \lookSecRef{sec:two-body-problem}.
